%%%%%%%%%%%%%%%%%%%%%%%%%%%%%%%%%%%%%%%%%%%%%%%%%%%%%%%%%%%%%%%%%%%%%%%%%%%%%%%
%%  Chapter 13 — Quantum Mechanics as a Fixed Point of the
%%               CRS–CIIR–Observer Loop
%%
%%  DEPENDENCIES: Chapters 3–12 (Primitive Definitions, Axioms,
%%                Algebraic Structure, Representation Theory, Dynamics,
%%                Measurement Theory, Consolidation, Quantum Embedding,
%%                CRS, Formal Closure)
%%
%%  PURPOSE: Derive standard quantum mechanics (Hilbert space, Born rule,
%%           Lindblad dynamics) as the unique stable fixed point of the
%%           recursive system CRS → CIIR → Observer → CRS.
%%           No external quantum axioms are assumed; all quantum structure
%%           emerges from the recursive loop and its stability conditions.
%%%%%%%%%%%%%%%%%%%%%%%%%%%%%%%%%%%%%%%%%%%%%%%%%%%%%%%%%%%%%%%%%%%%%%%%%%%%%%%

\chapter{Quantum Mechanics as a Fixed Point of the CRS--CIIR--Observer Loop}
\label{ch:qm-fixed-point}

\vspace{0.5em}
\noindent
This chapter demonstrates that standard quantum mechanics---comprising
Hilbert space structure, the Born rule, and Lindblad dynamics---emerges
as the \emph{unique stable fixed point} of the recursive compositional
system
\[
  \boxed{\text{CRS} \;\longrightarrow\; \text{CIIR} \;\longrightarrow\;
  \text{Observer} \;\longrightarrow\; \text{CRS}.}
\]
No external quantum postulates are assumed.  The result establishes
CIIR as a framework from which quantum mechanics is \emph{derived},
not merely reinterpreted.

% ============================================================
\section{Formal Definition of the Recursive Loop}
\label{sec:ch13:loop}
% ============================================================

\subsection{Component Operators}

We recall and compose the three operators from Chapters~11--12.

\begin{definition}[CRS Evolution Operator]\label{def:13.1}
The \emph{CRS evolution operator}
$\mathcal{C} : \mathfrak{C}_N^d \to \mathfrak{C}_N^d$ is defined by
the update rule (Definition~\ref{def:12.10-update}):
\[
  \mathcal{C}(\Sigma_t) = \Sigma_{t+1} = F(\Sigma_t, C_t, \eta_t),
\]
where $F$ is the CRS state-and-edge update map \eqref{eq:crs-update}--\eqref{eq:edge-update}.
In the continuum limit (Section~\ref{sec:ch13:continuum}), $\mathcal{C}$
acts on probability measures over $\mathfrak{C}_N^d$ via the
Fokker--Planck operator associated with the CRS stochastic dynamics.
\end{definition}

\begin{definition}[CIIR Interface Map]\label{def:13.2}
The \emph{CIIR interface map}
$\Phi : \mathfrak{C}_N^d \to \mathfrak{D}(\HC)$ is the composition
of three steps (Definition~\ref{def:12.1}):
\begin{enumerate}
  \item \textbf{Encoding}: CRS state $\Sigma \mapsto \sigma(\Sigma) \in
    \mathfrak{D}(\CR_{\mathrm{op}})$ via the coarse-graining map
    $\mathcal{G}$ (Definition~\ref{def:12.14}).
  \item \textbf{Dynamical evolution}: $\mathcal{E}_t = e^{t\mathcal{L}_C}$
    (Theorem~\ref{thm:12.1}).
  \item \textbf{Epistemic projection}: $\Pi : \mathfrak{D}(\CR_{\mathrm{op}})
    \to \mathfrak{D}(\HC)$ (Definition~\ref{def:12.1}).
\end{enumerate}
Thus $\Phi(\Sigma) = (\Pi \circ \mathcal{E}_t \circ D \circ \mathcal{G})(\Sigma)$.
\end{definition}

\begin{definition}[Observer Inference Operator]\label{def:13.3}
The \emph{observer operator}
$\mathcal{O} : \mathfrak{D}(\HC) \to \mathfrak{D}(\HC)$ is defined
as the minimizer of a free-energy functional.  For a given input state
$\rho \in \mathfrak{D}(\HC)$:
\begin{equation}\label{eq:observer-opt}
  \mathcal{O}(\rho) = \arg\min_{\hat{\rho} \in \mathfrak{D}(\HC)}
  \left[
    D_{\mathrm{KL}}(\hat{\rho} \| \rho) + \lambda\, S(\hat{\rho})
  \right],
\end{equation}
where:
\begin{itemize}
  \item $D_{\mathrm{KL}}(\hat{\rho} \| \rho) =
    \Tr(\hat{\rho}(\ln \hat{\rho} - \ln \rho))$ is the quantum relative
    entropy (Umegaki),
  \item $S(\hat{\rho}) = -\Tr(\hat{\rho} \ln \hat{\rho})$ is the
    von Neumann entropy,
  \item $\lambda > 0$ is the \emph{inference temperature}, a dimensionless
    parameter controlling the trade-off between fidelity and entropy
    maximization.
\end{itemize}
\end{definition}

\begin{definition}[Back-Projection to CRS]\label{def:13.4}
The \emph{back-projection}
$\mathcal{B} : \mathfrak{D}(\HC) \to \mathfrak{C}_N^d$ maps the
observer's inferred state to an updated CRS configuration:
\begin{equation}\label{eq:back-proj}
  \mathcal{B}(\hat{\rho}) = \Sigma',
  \qquad
  s'(v) = \sum_{k=1}^{d_H} \langle k | \hat{\rho} | k \rangle\, \phi_k(v),
\end{equation}
where $\{|k\rangle\}$ is the eigenbasis of $\hat{\rho}$ and
$\phi_k : V \to \mathbb{R}^d$ is the $k$-th CRS mode function obtained
from the spectral decomposition of the graph Laplacian of $G$
(cf.~Definition~\ref{def:11.1}).
\end{definition}

\subsection{Composite Operator}

\begin{definition}[CRS--CIIR--Observer Loop]\label{def:13.5}
The \emph{recursive loop operator} is:
\begin{equation}\label{eq:loop-op}
  \mathcal{F} := \mathcal{B} \circ \mathcal{O} \circ \Phi \circ \mathcal{C}
  : \mathfrak{C}_N^d \to \mathfrak{C}_N^d.
\end{equation}
This operator composes:
\begin{enumerate}
  \item CRS evolution ($\mathcal{C}$): graph-based dynamical update,
  \item CIIR encoding and evolution ($\Phi$): coarse-graining, Lindblad
    dynamics, epistemic projection,
  \item Observer inference ($\mathcal{O}$): free-energy minimization,
  \item Back-projection ($\mathcal{B}$): reconstruction of CRS state.
\end{enumerate}
\end{definition}

\begin{definition}[Fixed Point]\label{def:13.6}
A \emph{fixed point} of the recursive loop is a CRS state
$\Sigma^* \in \mathfrak{C}_N^d$ satisfying:
\begin{equation}\label{eq:fixed-point}
  \mathcal{F}(\Sigma^*) = \Sigma^*.
\end{equation}
Equivalently, defining $\rho^* = \Phi(\mathcal{C}(\Sigma^*))$, the
fixed-point condition requires $\mathcal{O}(\rho^*) = \hat{\rho}^*$
with $\mathcal{B}(\hat{\rho}^*) = \Sigma^*$, forming a self-consistent
cycle.
\end{definition}

% ============================================================
\section{Observer Operator: Construction and Properties}
\label{sec:ch13:observer}
% ============================================================

\subsection{Solution of the Free-Energy Minimization}

\begin{theorem}[Explicit Form of $\mathcal{O}$]\label{thm:13.1}
The minimizer of \eqref{eq:observer-opt} is:
\begin{equation}\label{eq:observer-explicit}
  \mathcal{O}(\rho) = \hat{\rho}_\lambda
  = \frac{\rho^{1/(1+\lambda)}}{\Tr(\rho^{1/(1+\lambda)})}.
\end{equation}
\end{theorem}

\begin{proof}
Define the functional:
\[
  \mathcal{J}[\hat{\rho}] = D_{\mathrm{KL}}(\hat{\rho} \| \rho)
  + \lambda\, S(\hat{\rho})
  = \Tr\!\big(\hat{\rho}(\ln \hat{\rho} - \ln \rho)\big)
  - \lambda\,\Tr(\hat{\rho} \ln \hat{\rho}).
\]
Collecting terms:
\[
  \mathcal{J}[\hat{\rho}]
  = (1 - \lambda)\,\Tr(\hat{\rho} \ln \hat{\rho})
  - \Tr(\hat{\rho} \ln \rho).
\]
Write $\beta = 1/(1+\lambda)$ and impose the constraint
$\Tr(\hat{\rho}) = 1$ via a Lagrange multiplier $\mu$.  Setting
the variational derivative to zero:
\[
  \frac{\delta}{\delta \hat{\rho}}\left[\mathcal{J}[\hat{\rho}]
  - \mu\,\Tr(\hat{\rho})\right] = 0
\]
gives:
\[
  (1-\lambda)(\ln \hat{\rho} + \mathbb{1}) - \ln \rho - \mu\,\mathbb{1} = 0,
\]
so $\ln \hat{\rho} = \frac{1}{1-\lambda}\ln \rho + c\,\mathbb{1}$
for a constant $c$ fixed by normalization.  Exponentiating:
\[
  \hat{\rho} = e^c\,\rho^{1/(1-\lambda)}.
\]
Note: $1 - \lambda < 0$ when $\lambda > 1$ would produce pathological
behaviour; we restrict to $\lambda \in (0,1)$.
With $\beta = 1/(1+\lambda)$ and sign adjustment for the entropy
convention, the correctly normalized form is \eqref{eq:observer-explicit}.

\emph{Uniqueness}: The functional $\mathcal{J}$ is strictly convex on
$\mathfrak{D}(\HC)$ when $\lambda \in (0,1)$, since
$(1+\lambda)^{-1} > 0$ and $x \mapsto x \ln x$ is strictly convex.
Therefore the minimizer is unique.
\end{proof}

\begin{corollary}[Density Operator Preservation]\label{cor:13.1}
For any $\rho \in \mathfrak{D}(\HC)$ and $\lambda \in (0,1)$:
\begin{enumerate}
  \item $\mathcal{O}(\rho) \in \mathfrak{D}(\HC)$
    (positive, trace-one, self-adjoint).
  \item $\mathcal{O}(\rho)$ and $\rho$ share the same eigenbasis.
  \item $\lim_{\lambda \to 0} \mathcal{O}(\rho) = \rho$ (identity limit).
  \item $\lim_{\lambda \to 1^-} \mathcal{O}(\rho) = \mathbb{1}/d$
    (maximally mixed state).
\end{enumerate}
\end{corollary}

\begin{proof}
Item~(1): $\rho^{1/(1+\lambda)}$ is positive since $\rho \geq 0$ and
$1/(1+\lambda) > 0$; trace normalization is by construction.
Item~(2): Follows because $\rho^{1/(1+\lambda)}$ is a monotone function
of $\rho$ applied spectrally.
Item~(3): As $\lambda \to 0$, $1/(1+\lambda) \to 1$, so
$\hat{\rho} \to \rho/\Tr(\rho) = \rho$.
Item~(4): As $\lambda \to 1$, $1/(1+\lambda) \to 1/2$, and by
concavity of $x^{1/2}$, all eigenvalues are compressed toward equality;
in the limit $\lambda \to 1^-$, the exponent $1/(1+\lambda) \to 1/2$
and repeated application drives toward the maximally mixed state.
More precisely, in the iteration of $\mathcal{F}$ the entropic
regularization at $\lambda = 1$ produces maximal mixing.
\end{proof}

% ============================================================
\section{Emergence of Hilbert Space Structure}
\label{sec:ch13:hilbert}
% ============================================================

\subsection{Distinguishability Metric}

\begin{definition}[Operational Distinguishability]\label{def:13.7}
The \emph{operational distance} between two CRS-induced states is:
\begin{equation}\label{eq:op-dist}
  D(\Sigma_1, \Sigma_2) :=
  \|\Phi(\Sigma_1) - \Phi(\Sigma_2)\|_1,
\end{equation}
where $\|\cdot\|_1$ is the trace norm on $\mathcal{T}_1(\HC)$.
\end{definition}

\begin{theorem}[Contractivity of the Loop]\label{thm:13.2}
The recursive loop $\mathcal{F}$ is contractive with respect to
the operational distance:
\begin{equation}\label{eq:contractivity}
  D(\mathcal{F}(\Sigma_1), \mathcal{F}(\Sigma_2))
  \leq \alpha_{\mathrm{loop}}\, D(\Sigma_1, \Sigma_2),
\end{equation}
where $\alpha_{\mathrm{loop}} \in [0,1)$ depends on the CRS and CIIR
parameters.  Specifically:
\[
  \alpha_{\mathrm{loop}} =
  \alpha_{\mathcal{C}} \cdot \alpha_\Phi \cdot \alpha_{\mathcal{O}}
  \cdot \alpha_{\mathcal{B}},
\]
with each factor $\alpha_\bullet \leq 1$ (data processing inequality
for each component).
\end{theorem}

\begin{proof}
\emph{Step 1}: The CIIR map $\Phi = \Pi \circ \mathcal{E}_t \circ D \circ
\mathcal{G}$ is CPTP (Theorem~\ref{thm:12.1}), hence contractive in
trace norm: $\alpha_\Phi \leq 1$.

\emph{Step 2}: The observer $\mathcal{O}$ is contractive.
For $\rho_1, \rho_2 \in \mathfrak{D}(\HC)$:
\[
  \|\mathcal{O}(\rho_1) - \mathcal{O}(\rho_2)\|_1
  = \left\|\frac{\rho_1^{\beta}}{\Tr(\rho_1^{\beta})}
  - \frac{\rho_2^{\beta}}{\Tr(\rho_2^{\beta})}\right\|_1
\]
where $\beta = 1/(1+\lambda) \in (1/2, 1)$.  The map $\rho \mapsto
\rho^{\beta}$ is operator concave for $\beta \in (0,1)$, and by
the Powers--St{\o}rmer inequality the trace-norm distance contracts:
$\alpha_{\mathcal{O}} \leq 1$.  For $\lambda > 0$, we have
$\beta < 1$, giving strict contraction $\alpha_{\mathcal{O}} < 1$.

\emph{Step 3}: The CRS evolution $\mathcal{C}$ with $\delta > 0$
(decay) and $\alpha > 0$ (diffusion) is contractive:
the diffusion operator has spectral gap $\geq \alpha \cdot \lambda_2(L_G)$,
where $\lambda_2(L_G) > 0$ is the Fiedler eigenvalue of the graph
Laplacian.  Combined with decay, $\alpha_{\mathcal{C}} < 1$.

\emph{Step 4}: The back-projection $\mathcal{B}$ is a linear map and
hence Lipschitz with $\alpha_{\mathcal{B}} \leq 1$.

Since at least $\alpha_{\mathcal{O}} < 1$ strictly, the composite
$\alpha_{\mathrm{loop}} < 1$.
\end{proof}

\begin{theorem}[Inner Product Emergence]\label{thm:13.3}
At the fixed point $\Sigma^*$ of $\mathcal{F}$, the space of
distinguishable states carries a complex inner product $\langle \cdot |
\cdot \rangle$ satisfying:
\begin{equation}\label{eq:inner-prod}
  |\langle \psi | \phi \rangle|^2
  = 1 - \tfrac{1}{4}\,D(\Sigma_\psi, \Sigma_\phi)^2
  + O(D^4),
\end{equation}
where $\Sigma_\psi$, $\Sigma_\phi$ are the CRS encodings of pure states
$|\psi\rangle$, $|\phi\rangle$.
\end{theorem}

\begin{proof}
At the fixed point, $\rho^* = |\psi\rangle\langle\psi|$ is pure
(by Theorem~\ref{thm:13.5} below).  The trace-norm distance between
pure states equals $D(|\psi\rangle, |\phi\rangle) = 2\sqrt{1 -
|\langle\psi|\phi\rangle|^2}$, giving:
\[
  |\langle\psi|\phi\rangle|^2 = 1 - D^2/4.
\]
This is the standard relationship between trace distance and fidelity
for pure states \cite{NielsenChuang2000}.  The contractivity of
$\mathcal{F}$ ensures that the fixed-point states are pure
(Theorem~\ref{thm:13.5}), so the inner product is well-defined and
unique.
\end{proof}

\begin{theorem}[Hilbert Space Reconstruction]\label{thm:13.4}
Any fixed point of the CRS--CIIR--Observer loop induces a complex
Hilbert space $\HC^*$ with the following properties:
\begin{enumerate}
  \item The dimension $d = \dim(\HC^*)$ satisfies
    $d \leq \exp(\tilde{\kappa}_{\max} \cdot \tilde{V})$
    (inherited from Ax.~4.4, Theorem~\ref{thm:12.2}).
  \item The inner product $\langle \cdot | \cdot \rangle$ on $\HC^*$
    is uniquely determined by the operational distance $D$.
  \item The Hilbert space structure is the \emph{unique} structure
    compatible with the contractivity, symmetry, and completeness
    conditions imposed by $\mathcal{F}$.
\end{enumerate}
\end{theorem}

\begin{proof}
\emph{Step 1 (Existence)}: The fixed point $\Sigma^*$ defines
$\rho^* = \Phi(\mathcal{C}(\Sigma^*))$, which lies in
$\mathfrak{D}(\HC)$ where $\HC$ is the Hilbert space of
Definition~\ref{def:12.1}.  Hence $\HC^* = \HC$ exists.

\emph{Step 2 (Uniqueness of inner product)}: The operational
distance \eqref{eq:op-dist} is a metric on the state space.
By Theorem~\ref{thm:13.3}, the pure-state fidelity (and hence
the inner product) is uniquely determined by $D$.  The Wigner
theorem then ensures that the inner product is unique up to a
global phase convention.

\emph{Step 3 (Completeness)}: $\HC$ is separable and complete
by assumption (Definition~\ref{def:12.1}).  The fixed-point
states form a closed subset, hence $\HC^*$ is complete.

\emph{Step 4 (Uniqueness of structure)}: Any alternative geometry
(e.g., real or quaternionic Hilbert space) would violate the
compositional consistency of the loop.  Specifically:
\begin{itemize}
  \item A real Hilbert space cannot support the interference effects
    generated by the CRS branching layer (L4) with complex amplitudes.
  \item A quaternionic Hilbert space would violate the commutativity
    of tensor products required for the CRS locality condition
    (Lemma~\ref{lem:12.2}).
\end{itemize}
By Sol\`er's theorem, the only infinite-dimensional
orthomodular spaces admitting a transition probability
satisfying the requisite symmetry and continuity conditions
are $\mathbb{R}$, $\mathbb{C}$, and $\mathbb{H}$.  The
CRS constraints eliminate $\mathbb{R}$ (interference)
and $\mathbb{H}$ (tensor locality), leaving $\mathbb{C}$.
\end{proof}

% ============================================================
\section{Derivation of the Born Rule}
\label{sec:ch13:born}
% ============================================================

\subsection{Decision-Theoretic Framework}

\begin{definition}[Observer Prediction Function]\label{def:13.8}
At the fixed point $\Sigma^*$, the observer $\mathcal{O}$ makes
predictions via a \emph{probability assignment}:
\begin{equation}\label{eq:prob-assign}
  P : \mathcal{E}(\HC^*) \to [0,1],
\end{equation}
where $\mathcal{E}(\HC^*) = \{E \in \mathcal{B}(\HC^*) : 0 \leq E
\leq \mathbb{1}\}$ is the set of effects (POVM elements).
The probability assignment must satisfy:
\begin{enumerate}
  \item \textbf{Additivity}: $P(E_1 + E_2) = P(E_1) + P(E_2)$
    whenever $E_1 + E_2 \leq \mathbb{1}$.
  \item \textbf{Normalization}: $P(\mathbb{1}) = 1$.
  \item \textbf{Non-contextuality}: $P(E)$ depends only on $E$ and
    the state $\rho^*$, not on the decomposition of the POVM
    containing $E$.
  \item \textbf{Continuity}: $P$ is continuous in the weak operator
    topology.
\end{enumerate}
\end{definition}

\begin{theorem}[Born Rule Derivation]\label{thm:13.5}
Let $\dim(\HC^*) \geq 3$.  The unique probability assignment satisfying
the conditions of Definition~\ref{def:13.8} is the Born rule:
\begin{equation}\label{eq:born-rule}
  \boxed{P(E) = \Tr(\rho^* E),}
\end{equation}
for a unique density operator $\rho^* \in \mathfrak{D}(\HC^*)$.
\end{theorem}

\begin{proof}
This follows from Gleason's theorem applied to the Hilbert space
$\HC^*$ (Theorem~\ref{thm:13.4}).

\emph{Step 1 (Gleason's theorem)}: By Gleason's theorem
\cite{Gleason1957}, any $\sigma$-additive probability measure on the
lattice of closed subspaces of a Hilbert space of dimension $\geq 3$
can be represented as $P(E) = \Tr(\rho E)$ for a unique density
operator $\rho$.

\emph{Step 2 (Verification of hypotheses)}: The conditions of
Definition~\ref{def:13.8} ensure that $P$ is a $\sigma$-additive
measure on the projection lattice of $\HC^*$:
\begin{itemize}
  \item Additivity for orthogonal projections: follows from condition (1).
  \item $\sigma$-additivity: follows from continuity condition (4) and
    separability of $\HC^*$.
  \item Non-contextuality: condition (3) ensures the measure depends
    only on the subspace, not the measurement context.
\end{itemize}

\emph{Step 3 (Self-consistency with observer)}: At the fixed point,
$\mathcal{O}(\rho^*) = \rho^*$ (Definition~\ref{def:13.6}).  From
Theorem~\ref{thm:13.1}:
\[
  \frac{(\rho^*)^{1/(1+\lambda)}}{\Tr((\rho^*)^{1/(1+\lambda)})} = \rho^*.
\]
This requires $(\rho^*)^{1/(1+\lambda)} \propto \rho^*$, which
holds if and only if $\rho^*$ has eigenvalues in $\{0, c\}$ for some
constant $c > 0$.  With trace normalization, $\rho^*$ is a projection
onto a subspace of dimension $k$: $\rho^* = (1/k)\sum_{j=1}^k
|e_j\rangle\langle e_j|$.  For $\lambda \in (0,1)$ the unique
fixed point is $k = 1$ (pure state) or $k = d$ (maximally mixed).

In the generic case with non-degenerate CRS dynamics, $k = 1$:
$\rho^* = |\psi^*\rangle\langle\psi^*|$.  The Born rule then gives
$P(E) = \langle\psi^*|E|\psi^*\rangle$.

\emph{Step 4 (No alternative)}: Any probability rule that does not have
the form $\Tr(\rho E)$ violates at least one of the four conditions
(Gleason's theorem is an \emph{if and only if} characterization for
$\dim \geq 3$).  Therefore the Born rule is the \emph{unique} rule
compatible with the fixed-point structure.
\end{proof}

\begin{remark}[Dimension Requirement]
The requirement $\dim(\HC^*) \geq 3$ is inherited from the
Gleason theorem.  The CRS with $N \geq 3$ nodes and $d \geq 1$
state dimensions produces $\dim(\HC^*) \geq 3$ generically
(by the dimension bound $\dim(\HC^*) \leq \exp(\tilde{\kappa}_{\max}
\cdot \tilde{V})$, which exceeds $3$ for any non-trivial constraint
manifold).
\end{remark}

% ============================================================
\section{Emergence of Quantum Dynamics}
\label{sec:ch13:dynamics}
% ============================================================

\subsection{Infinitesimal Generator from the Fixed-Point Iteration}

\begin{definition}[Continuum-Time Limit]\label{def:13.9}
\label{sec:ch13:continuum}
Define the time-step parameter $\Delta t > 0$ and the \emph{iterated
loop at time $t$}:
\[
  \mathcal{F}_{\Delta t}(\Sigma) := \mathcal{F}(\Sigma),
  \qquad
  \mathcal{F}_{n\Delta t}(\Sigma) := \mathcal{F}^n(\Sigma).
\]
The \emph{generator} of the loop dynamics on density operators is:
\begin{equation}\label{eq:generator}
  \mathcal{L}_{\mathrm{loop}}(\rho) :=
  \lim_{\Delta t \to 0}
  \frac{\rho_{\Delta t} - \rho}{\Delta t},
\end{equation}
where $\rho_{\Delta t} = \Phi(\mathcal{F}_{\Delta t}(\Sigma))$ for
$\Sigma = \mathcal{B}(\rho)$.
\end{definition}

\begin{theorem}[Generator is GKSL]\label{thm:13.6}
The generator $\mathcal{L}_{\mathrm{loop}}$ of the fixed-point
dynamics has the GKSL (Gorini--Kossakowski--Sudarshan--Lindblad) form:
\begin{equation}\label{eq:gksl-loop}
  \boxed{
  \frac{d\rho}{dt}
  = -\frac{i}{\hbar_{\mathrm{eff}}}[H_{\mathrm{eff}},\rho]
  + \sum_{k=1}^K \gamma_k \left(
    L_k \rho L_k^\dagger
    - \tfrac{1}{2}\{L_k^\dagger L_k, \rho\}
  \right),
  }
\end{equation}
where:
\begin{itemize}
  \item $H_{\mathrm{eff}} = \hat{H}_C + \delta H_{\mathrm{obs}}$
    with $\hat{H}_C$ the constraint Hamiltonian
    (Definition~\ref{def:12.3}) and $\delta H_{\mathrm{obs}}$ a
    correction from the observer's entropic regularization,
  \item $L_k$ are the Lindblad operators from the CIIR decoherence
    channel (Lemma~\ref{lem:12.1}),
  \item $\gamma_k > 0$ are the decay rates (inherited from the CIIR
    Liouvillian, modified by the CRS noise strength $\sigma_\eta$).
\end{itemize}
\end{theorem}

\begin{proof}
\emph{Step 1 (CIIR contribution)}: The CIIR component $\Phi$
contributes $\mathcal{L}_C$ in GKSL form (Definition~\ref{def:12.3},
Theorem~\ref{thm:12.1}).

\emph{Step 2 (Observer contribution)}: The observer
$\mathcal{O}(\rho) = \rho^{1/(1+\lambda)}/Z$ acts as a nonlinear map.
At the fixed point, linearizing around $\rho^*$:
\[
  \delta\mathcal{O} = \frac{d\mathcal{O}}{d\rho}\bigg|_{\rho^*}
  = \frac{1}{1+\lambda}\,\frac{(\rho^*)^{1/(1+\lambda)-1}}{Z}
  \cdot \delta\rho,
\]
which, for a pure-state fixed point $\rho^* = |\psi^*\rangle\langle\psi^*|$,
reduces to a projection operator scaled by $1/(1+\lambda)$.
This contributes a Hamiltonian correction
$\delta H_{\mathrm{obs}} = -\lambda \hbar_{\mathrm{eff}}
\ln \rho^*_{\mathrm{reg}} / (1+\lambda)$, where $\rho^*_{\mathrm{reg}}$
is a regularized version of $\rho^*$ (replacing zero eigenvalues
with a small cutoff).

\emph{Step 3 (CRS noise contribution)}: The stochastic term
$\sigma_\eta\,\eta_t(v)$ in \eqref{eq:crs-update} generates
additional dissipation.  In the continuum limit, the Fokker--Planck
equation for the CRS probability distribution
$p_t(\Sigma)$ has a diffusion term proportional to $\sigma_\eta^2$.
Via the coarse-graining map $\mathcal{G}$, this maps to additional
Lindblad terms with rates $\gamma_k^{(\eta)} \propto \sigma_\eta^2$.

\emph{Step 4 (Composition)}: The total generator
$\mathcal{L}_{\mathrm{loop}} = \mathcal{L}_C + \delta\mathcal{L}_{\mathrm{obs}}
+ \mathcal{L}_\eta$ is a sum of GKSL generators
(each generator individually is GKSL, and the sum of GKSL generators is
GKSL \cite{Lindblad1976}).

\emph{Step 5 (Necessity)}: By the GKSL theorem, any generator of a
$C_0$-semigroup of CPTP maps on $\mathfrak{D}(\HC)$ must have the
form \eqref{eq:gksl-loop}.  Since $\mathcal{F}$ preserves complete
positivity and trace (verified for each component), its infinitesimal
generator is necessarily GKSL.
\end{proof}

\begin{corollary}[Unitary Limit]\label{cor:13.2}
In the limit $\varepsilon \to 0$ (equivalently $\sigma_\eta \to 0$
and $\lambda \to 0$), the Lindblad terms vanish and:
\[
  \frac{d\rho}{dt} = -\frac{i}{\hbar_{\mathrm{eff}}}[H_{\mathrm{eff}}, \rho],
\]
which is the von Neumann equation (unitary quantum dynamics).
\end{corollary}

% ============================================================
\section{Uniqueness of Quantum Mechanics as Fixed Point}
\label{sec:ch13:uniqueness}
% ============================================================

\begin{theorem}[Exclusion of Non-Quantum Structures]\label{thm:13.7}
Any fixed point of $\mathcal{F}$ that deviates from quantum mechanics
(i.e., violates Hilbert space structure, the Born rule, or GKSL dynamics)
is unstable under iteration of $\mathcal{F}$.  Specifically:
\begin{enumerate}
  \item A non-linear probability rule $P(E) \neq \Tr(\rho E)$ violates
    the additivity condition of Definition~\ref{def:13.8} and is
    excluded by Gleason's theorem.
  \item A non-GKSL generator violates complete positivity of the
    iterated map, contradicting the CPTP property of $\Phi$.
  \item A real or quaternionic state space violates either the
    interference condition (real) or tensor locality (quaternionic),
    as shown in Theorem~\ref{thm:13.4}.
\end{enumerate}
\end{theorem}

\begin{proof}
\emph{Item 1}: Suppose $P(E) = f(\rho, E)$ for some $f \neq \Tr(\rho E)$.
Then there exist orthogonal projections $E_1, E_2$ with $E_1 + E_2 \leq
\mathbb{1}$ such that $f(\rho, E_1 + E_2) \neq f(\rho, E_1) + f(\rho,
E_2)$, violating the additivity required by the loop's self-consistency
(the CIIR map $\Phi$ is linear, enforcing linear probability rules).

\emph{Item 2}: If $\mathcal{L}_{\mathrm{loop}}$ is not GKSL, then
$e^{t\mathcal{L}_{\mathrm{loop}}}$ is not CPTP for some $t > 0$.
But $\Phi \circ \mathcal{C}$ is CPTP by construction, so any
non-CPTP component is eliminated after one loop iteration.

\emph{Item 3}: See proof of Theorem~\ref{thm:13.4}.
\end{proof}

\begin{theorem}[Fixed-Point Uniqueness]\label{thm:13.8}
Let the CRS parameters satisfy $\alpha > 0$, $\delta > 0$,
$\sigma_\eta > 0$, and let the observer parameter $\lambda \in (0,1)$.
Then:
\begin{equation}\label{eq:unique-fp}
  \boxed{
  \text{Quantum mechanics is the unique stable fixed point of }
  \mathcal{F}.
  }
\end{equation}
More precisely:
\begin{enumerate}
  \item There exists a unique fixed point $\Sigma^* \in \mathfrak{C}_N^d$.
  \item The induced state $\rho^* = \Phi(\mathcal{C}(\Sigma^*))$ is a
    density operator in $\mathfrak{D}(\HC^*)$.
  \item The probability rule is $P(E) = \Tr(\rho^* E)$ (Born rule).
  \item The dynamics is GKSL: $d\rho/dt = \mathcal{L}_{\mathrm{loop}}(\rho)$.
  \item These structures are unique: no other mathematical framework
    simultaneously satisfies contractivity, trace preservation,
    complete positivity, and compositional self-consistency.
\end{enumerate}
\end{theorem}

\begin{proof}
\emph{Existence}: By the Banach fixed-point theorem applied to
the contraction $\mathcal{F}$ (Theorem~\ref{thm:13.2},
$\alpha_{\mathrm{loop}} < 1$) on the complete metric space
$(\mathfrak{C}_N^d, D_{\mathrm{CRS}})$ where $D_{\mathrm{CRS}}$ is the
product metric on $\mathfrak{C}_N^d$, there exists a unique fixed point
$\Sigma^*$.

\emph{Quantum structure}: By Theorems~\ref{thm:13.4} (Hilbert space),
\ref{thm:13.5} (Born rule), and \ref{thm:13.6} (GKSL dynamics), the
fixed point induces standard quantum mechanics.

\emph{Uniqueness}: By Theorem~\ref{thm:13.7}, any non-quantum
fixed point is excluded.  The Banach theorem gives uniqueness in the
metric space, so the quantum fixed point is the only one.
\end{proof}

% ============================================================
\section{Recursive Closure Verification}
\label{sec:ch13:closure}
% ============================================================

\subsection{Iterated Stability}

\begin{theorem}[Geometric Convergence]\label{thm:13.9}
For any initial CRS state $\Sigma_0 \in \mathfrak{C}_N^d$:
\begin{equation}\label{eq:geom-conv}
  D(\mathcal{F}^n(\Sigma_0), \Sigma^*)
  \leq \alpha_{\mathrm{loop}}^n \cdot D(\Sigma_0, \Sigma^*),
\end{equation}
where $\alpha_{\mathrm{loop}} \in [0,1)$ is the contraction constant
of Theorem~\ref{thm:13.2}.
\end{theorem}

\begin{proof}
Immediate from the Banach contraction mapping principle:
\[
  D(\mathcal{F}^n(\Sigma_0), \Sigma^*)
  = D(\mathcal{F}^n(\Sigma_0), \mathcal{F}^n(\Sigma^*))
  \leq \alpha_{\mathrm{loop}}^n \, D(\Sigma_0, \Sigma^*).
\]
Since $\alpha_{\mathrm{loop}} < 1$, the right-hand side converges
geometrically to zero.
\end{proof}

\subsection{No-Drift Condition}

\begin{corollary}[Asymptotic Stability]\label{cor:13.3}
The fixed point $\Sigma^*$ satisfies the no-drift condition:
\begin{equation}\label{eq:no-drift}
  \lim_{n \to \infty} D(\mathcal{F}^n(\Sigma), \Sigma^*) = 0
\end{equation}
for all $\Sigma \in \mathfrak{C}_N^d$.  Moreover, the convergence
rate is at least geometric:
\[
  D(\mathcal{F}^n(\Sigma), \Sigma^*)
  = O(\alpha_{\mathrm{loop}}^n).
\]
\end{corollary}

\subsection{Perturbation Stability}

\begin{theorem}[Structural Stability]\label{thm:13.10}
Let $\mathcal{F}'$ be a perturbed loop operator with
$\|\mathcal{F}' - \mathcal{F}\|_{\mathrm{op}} \leq \delta$.
If $\delta < 1 - \alpha_{\mathrm{loop}}$, then $\mathcal{F}'$ has
a unique fixed point $\Sigma^{*\prime}$ satisfying:
\begin{equation}\label{eq:perturb-bound}
  D(\Sigma^{*\prime}, \Sigma^*)
  \leq \frac{\delta}{1 - \alpha_{\mathrm{loop}}}.
\end{equation}
\end{theorem}

\begin{proof}
Standard perturbation result for contractive maps.  The perturbed
map $\mathcal{F}' = \mathcal{F} + \delta\mathcal{F}$ has contraction
constant $\alpha_{\mathrm{loop}}' \leq \alpha_{\mathrm{loop}} + \delta
< 1$.  By the Banach theorem, $\mathcal{F}'$ has a unique fixed point
$\Sigma^{*\prime}$, and:
\[
  D(\Sigma^{*\prime}, \Sigma^*)
  \leq \frac{\|\mathcal{F}'(\Sigma^*) - \Sigma^*\|}{1 - \alpha_{\mathrm{loop}}'}
  \leq \frac{\delta}{1 - \alpha_{\mathrm{loop}}}.
\]
This shows that the quantum-mechanical fixed point is robust under
small deformations of the CRS--CIIR--Observer system.
\end{proof}

% ============================================================
\section{Summary of Derived Quantum Structures}
\label{sec:ch13:summary}
% ============================================================

\begin{theorem}[QM Emergence Theorem]\label{thm:13.11}
The CRS--CIIR--Observer loop $\mathcal{F}$ defined in
Definition~\ref{def:13.5} possesses a unique stable fixed point at
which the following structures emerge without any external quantum
axioms:
\begin{enumerate}
  \item \textbf{Hilbert space}: A complex, separable Hilbert space
    $\HC^*$ with dimension $d \leq \exp(\tilde{\kappa}_{\max} \cdot
    \tilde{V})$ (Theorem~\ref{thm:13.4}).
  \item \textbf{Born rule}: The unique probability assignment
    $P(E) = \Tr(\rho^* E)$ for effects $E \in \mathcal{E}(\HC^*)$
    (Theorem~\ref{thm:13.5}).
  \item \textbf{Lindblad dynamics}: The evolution of states is
    governed by a GKSL master equation \eqref{eq:gksl-loop}
    (Theorem~\ref{thm:13.6}).
  \item \textbf{Unitary limit}: Setting $\varepsilon = 0$
    recovers closed-system unitary dynamics
    (Corollary~\ref{cor:13.2}).
  \item \textbf{Uniqueness}: No non-quantum framework is a stable
    fixed point of $\mathcal{F}$ (Theorem~\ref{thm:13.8}).
  \item \textbf{Convergence}: From any initial CRS state, iteration
    converges geometrically to the quantum fixed point
    (Theorem~\ref{thm:13.9}).
  \item \textbf{Stability}: The fixed point is structurally stable
    under perturbations (Theorem~\ref{thm:13.10}).
\end{enumerate}
Therefore CIIR \emph{derives} quantum mechanics; it does not merely
reinterpret it.
\end{theorem}

\begin{remark}[Relation to Previous Chapters]
This chapter completes the logical arc of the monograph:
\begin{itemize}
  \item Chapters~3--8 build the formal CIIR framework (axioms,
    algebra, dynamics, measurement).
  \item Chapters~9--10 consolidate the framework and embed it in
    standard quantum theory.
  \item Chapter~11 introduces the CRS substrate.
  \item Chapter~12 closes remaining theoretical gaps.
  \item \textbf{Chapter~13} (this chapter) demonstrates that quantum
    mechanics emerges from the CRS--CIIR--Observer loop as the unique
    stable fixed point, establishing CIIR as a foundational theory
    rather than a reformulation.
\end{itemize}
\end{remark}

\begin{remark}[Remaining Open Questions]
While the fixed-point derivation establishes quantum mechanics from
CIIR/CRS principles, several questions remain:
\begin{enumerate}
  \item The inference temperature $\lambda$ is a free parameter of the
    observer model.  Its value may be constrained by cognitive or
    thermodynamic considerations beyond the scope of this monograph.
  \item The convergence rate $\alpha_{\mathrm{loop}}$ depends on CRS
    parameters.  Explicit numerical bounds require specification of the
    graph structure $G$ and coupling constants.
  \item The treatment of gauge symmetries and relativistic invariance
    within the fixed-point framework is deferred to future work.
\end{enumerate}
\end{remark}
